% !TEX program = pdflatex
\documentclass[11pt,a4paper]{article}

% =========================================================================
% PLAN DE CALIDAD
% Proyecto: CasaScore (AlquilaPuno)
% Gestión de Proyectos TICS (EST315) -- Semestre 2026-II -- Team Sobrevivientes
% =========================================================================

% --- PAQUETES DE CONFIGURACIÓN ---
\usepackage[utf8]{inputenc}
\usepackage[T1]{fontenc}
\usepackage[spanish,es-tabla]{babel}
% Compatibilidad con babel-spanish reciente (flechas de TikZ y signo de porcentaje)
\AtBeginDocument{\shorthandoff{<>}\spanishplainpercent}
\usepackage{geometry}
\geometry{
    top=3.6cm,
    bottom=2.3cm,
    left=2.2cm,
    right=2.2cm,
    % PARTE MODIFICABLE: headheight debe ser mayor que la altura del logo más alto (2.56cm = 73pt).
    headheight=80pt,
    headsep=15pt,
    footskip=20pt
}

\usepackage{graphicx}
\usepackage{fancyhdr}
\usepackage[table]{xcolor}
\usepackage{titlesec}
\usepackage{booktabs}
\usepackage{tabularx}
\usepackage{array}
\usepackage{amsmath}
\usepackage{enumitem}
\usepackage{microtype}
\usepackage{amssymb}
% --- RECURSOS VISUALES: GRÁFICOS Y CÓDIGO DE COLORES ---
\usepackage{tikz}
\usetikzlibrary{positioning,calc}
\definecolor{GrisFila}{RGB}{246, 248, 251}
\definecolor{CuadCerca}{RGB}{150, 30, 30}
\definecolor{CuadSatisf}{RGB}{180, 130, 20}
\definecolor{CuadInfo}{RGB}{10, 45, 90}
\definecolor{CuadMonit}{RGB}{110, 115, 125}
\newcommand{\etiqueta}[2]{{\setlength{\fboxsep}{2pt}\colorbox{#1}{\textcolor{white}{\scriptsize\bfseries #2}}}}
\newcommand{\cerca}{\etiqueta{CuadCerca}{Gestionar de cerca}}
\newcommand{\satisf}{\etiqueta{CuadSatisf}{Mantener satisfecho}}
\newcommand{\informado}{\etiqueta{CuadInfo}{Mantener informado}}
\newcommand{\monit}{\etiqueta{CuadMonit}{Monitorear}}
\usepackage{needspace}
\usepackage{longtable}
\definecolor{Rama1}{RGB}{10, 45, 90}
\definecolor{Rama2}{RGB}{40, 95, 150}
\definecolor{Rama3}{RGB}{180, 130, 20}
\definecolor{Rama4}{RGB}{30, 110, 60}
\definecolor{Rama5}{RGB}{110, 60, 140}
\definecolor{RaciR}{RGB}{150, 30, 30}
\definecolor{RaciA}{RGB}{10, 45, 90}
\definecolor{RaciC}{RGB}{180, 130, 20}
\definecolor{RaciI}{RGB}{140, 145, 155}
\definecolor{NivAlto}{RGB}{30, 110, 60}
\definecolor{NivMedio}{RGB}{180, 130, 20}
\definecolor{NivBajo}{RGB}{150, 30, 30}
\usetikzlibrary{patterns}
\definecolor{SevBaja}{RGB}{30, 110, 60}
\definecolor{SevMedia}{RGB}{180, 130, 20}
\definecolor{SevAlta}{RGB}{150, 30, 30}
\usetikzlibrary{arrows.meta,shapes.geometric,fit,backgrounds}
\usepackage{listings}
\lstset{basicstyle=\ttfamily\scriptsize, frame=single, rulecolor=\color{AzulUNA!40}, backgroundcolor=\color{GrisFila}, columns=fullflexible, keepspaces=true, breaklines=true, xleftmargin=4pt, xrightmargin=4pt, literate={á}{{\'a}}1 {é}{{\'e}}1 {í}{{\'i}}1 {ó}{{\'o}}1 {ú}{{\'u}}1 {ñ}{{\~n}}1 {Á}{{\'A}}1}

\usepackage{hyperref}

% Logos: funciona a 2 o 3 niveles de profundidad dentro de docs/latex/
\graphicspath{{../../assets/img/}{../../../assets/img/}}

% --- PALETA INSTITUCIONAL FIEI / UNA PUNO ---
\definecolor{AzulUNA}{RGB}{10, 45, 90}
\definecolor{DoradoFIEI}{RGB}{180, 130, 20}
\definecolor{GrisOscuro}{RGB}{50, 50, 50}
\definecolor{AzulClaro}{RGB}{228, 236, 246}

% --- CONFIGURACIÓN DE LISTAS ---
\setlist{topsep=2pt, itemsep=2pt, parsep=1pt}

% --- COLUMNAS DE TABLA ---
\newcolumntype{L}{>{\raggedright\arraybackslash}X}
\newcolumntype{P}[1]{>{\raggedright\arraybackslash}p{#1}}
\newcolumntype{C}[1]{>{\centering\arraybackslash}p{#1}}
\newcolumntype{Q}[1]{>{\raggedleft\arraybackslash}p{#1}}
\newcolumntype{R}{>{\raggedleft\arraybackslash}X}
\renewcommand{\arraystretch}{1.25}

% --- FORMATO DE ENLACES ---
\hypersetup{
    colorlinks=true,
    linkcolor=AzulUNA,
    urlcolor=AzulUNA,
    citecolor=AzulUNA,
    pdftitle={Plan de Calidad -- CasaScore (AlquilaPuno)},
    pdfauthor={Team Sobrevivientes}
}

% --- FORMATO DE TÍTULOS ---
\titleformat{\section}{\color{AzulUNA}\normalfont\large\bfseries}{\thesection.}{0.7em}{}[\color{AzulUNA}\titlerule]
\titleformat{\subsection}{\color{AzulUNA}\normalfont\normalsize\bfseries}{\thesubsection.}{0.5em}{}
\titlespacing*{\section}{0pt}{9pt}{4pt}
\titlespacing*{\subsection}{0pt}{6pt}{2pt}

% --- ENCABEZADO Y PIE DE PÁGINA (fancyhdr) ---
\pagestyle{fancy}
\fancyhf{}

% ===== TAMAÑO DE LOS LOGOS (PARTE MODIFICABLE) =====
% logo_fiei.png  : height=1.45cm (lado izquierdo)
% logo_equipo.png: height=2.56cm (lado derecho; 25 % más grande porque la imagen tiene más borde blanco)
% Para cambiar el tamaño, modifique solo el valor height=... de cada \includegraphics.
% Si el logo más alto supera 2.6cm, aumente también headheight y top en la configuración de página.
\fancyhead[L]{%
    \includegraphics[height=1.45cm,keepaspectratio]{logo_fiei.png}%
}
\fancyhead[C]{%
    \scriptsize \textbf{\textcolor{AzulUNA}{UNIVERSIDAD NACIONAL DEL ALTIPLANO DE PUNO}}\\[2pt]
    \tiny \textcolor{GrisOscuro}{FACULTAD DE INGENIERÍA ESTADÍSTICA E INFORMÁTICA}\\[2pt]
    \scriptsize \textbf{\textcolor{AzulUNA}{Gestión de Proyectos TICS (EST315) \textbar\ Team Sobrevivientes}}\\[2pt]
    \footnotesize \textbf{\textcolor{DoradoFIEI}{PLAN DE CALIDAD}}%
}
\fancyhead[R]{%
    \includegraphics[height=2.56cm,keepaspectratio]{logo_equipo.png}%
}

\renewcommand{\headrulewidth}{0pt}
\renewcommand{\headrule}{}

\fancyfoot[L]{\scriptsize \textcolor{GrisOscuro}{CasaScore (AlquilaPuno)}}
\fancyfoot[C]{\small \thepage}
\fancyfoot[R]{\scriptsize \textcolor{GrisOscuro}{Puno, Perú -- Semestre 2026-II}}
\renewcommand{\footrulewidth}{0.4pt}

\begin{document}

\begin{center}
    {\Large \textbf{\textcolor{AzulUNA}{PLAN DE CALIDAD}}}\\[3pt]
    {\normalsize \textbf{CasaScore (AlquilaPuno)}}\\[4pt]
    \footnotesize \textbf{Fecha:} martes 22/09/2026 (semana 5) \quad \textbar\quad \textbf{Versión:} 1.0 \quad \textbar\quad \textbf{Fase:} Planificación
\end{center}

\vspace{0.1cm}

\section{Datos Generales}
\noindent\begin{tabularx}{\textwidth}{@{}P{4.2cm} L@{}}
\toprule
\textbf{Proyecto:} & CasaScore (aplicación AlquilaPuno) \\
\textbf{Responsable:} & Eliseo Tarqui Ajahuana -- Especialista Técnico / Arquitectura \\
\textbf{Proceso PM4DEV:} & Gestión de la calidad \\
\textbf{Insumos:} & Especificación de requerimientos, Acta de Constitución, EDT y cronograma \\
\bottomrule
\end{tabularx}

\section{Política de Calidad}
CasaScore busca que el arrendatario confíe en la recomendación y que el propietario publique con facilidad. La calidad se mide por el cumplimiento de los requisitos y por la satisfacción de ambos clientes: \textbf{80\,\% o más califica la app con 4 o 5 sobre 5} en el piloto. La calidad se planifica, se construye en cada sprint y se verifica antes de cada entrega; no se deja para el final.

\section{Estándares}
\begin{itemize}[leftmargin=*]
    \item \textbf{Producto:} modelo de calidad ISO/IEC 25010:2023, aplicado a las características que importan al proyecto.
    \item \textbf{Código:} TypeScript en modo estricto, motor de puntaje separado de las pantallas, nombres en inglés y mensajes de commit con el formato \texttt{feat(casascore-...)}.
    \item \textbf{Documentos:} plantilla \LaTeX\ del equipo, guía PM4DEV y redacción formal.
    \item \textbf{Datos personales:} Ley N.\textsuperscript{o} 29733, Ley de Protección de Datos Personales.
\end{itemize}

\Needspace{12\baselineskip}
\section{Características de Calidad del Producto}
{\footnotesize
\setlength{\tabcolsep}{3pt}
\rowcolors{2}{white}{GrisFila}
\noindent\begin{tabularx}{\textwidth}{@{}P{2.35cm} P{2.45cm} P{1.6cm} L C{1.45cm} P{2.25cm}@{}}
\toprule
\rowcolor{AzulClaro}
\textbf{Característica} & \textbf{Subcaracte\-rística} & \textbf{Requisito} & \textbf{Métrica} & \textbf{Meta} & \textbf{Verificación} \\
\midrule
\textbf{Adecuación funcional} & Completitud y corrección & RF-01 a RF-09; RNF-06 & Requisitos aceptados; puntajes con desglose correcto & 100\,\% & Pruebas de aceptación; CP-07 y CP-11 \\
\textbf{Eficiencia de desempeño} & Comportamiento temporal & RNF-01 & Tiempo para calcular y mostrar el ranking con 500 cuartos & $\le$ 250 ms & CP-12 \\
\textbf{Capacidad de interacción (usabilidad)} & Operabilidad y facilidad de aprendizaje & RNF-05 & Tiempo para publicar un cuarto sin ayuda; satisfacción en el piloto & $\le$ 10 min; $\ge$ 80\,\% & CP-09 y encuesta del piloto \\
\textbf{Fiabilidad} & Ausencia de fallos & O2 del Acta & Errores críticos abiertos al entregar el MVP & 0 & Registro de defectos \\
\textbf{Seguridad} & Confidencialidad e integridad & RNF-03 y RNF-04 & Accesos indebidos entre roles en las pruebas & 0 & CP-08 y CP-13 \\
\textbf{Mantenibilidad} & Modularidad y capacidad de prueba & RNF-08 & Cobertura de pruebas automáticas del motor de puntaje & $\ge$ 80\,\% & Reporte de cobertura de Jest \\
\textbf{Flexibilidad (portabilidad)} & Adaptabilidad & RNF-02 & Funciona en Android 10 y en una versión reciente & 2 de 2 equipos & CP-14 \\
\bottomrule
\end{tabularx}
}

\Needspace{14\baselineskip}
\section{Aseguramiento de la Calidad}
El aseguramiento revisa los procesos para prevenir errores antes de que lleguen al producto.

\vspace{0.1cm}
{\small
\rowcolors{2}{white}{GrisFila}
\noindent\begin{tabularx}{\textwidth}{@{}P{3.4cm} L P{2.5cm} P{2.3cm}@{}}
\toprule
\rowcolor{AzulClaro}
\textbf{Actividad} & \textbf{Qué se revisa} & \textbf{Frecuencia} & \textbf{Responsable} \\
\midrule
Revisión de código & Cada Pull Request lo revisa un integrante distinto al autor antes de unirlo a la rama principal. & Cada Pull Request & Eliseo Tarqui Ajahuana \\
Revisión de documentos & Otro integrante revisa contenido, cifras y formato antes de entregar. & Cada documento & Viviana M. Quisbert Quispe \\
Definición de terminado & Se verifica la lista de la sección 8 antes de cerrar un Issue. & Cada Issue & Responsable del Issue \\
Retrospectiva & Qué funcionó, qué falló y qué se cambia en el siguiente sprint. & Viernes & Clyde N. Paricahua Pari \\
Auditoría interna & Cumplimiento de este plan y de los criterios de aceptación. & Semana 13 & Viviana M. Quisbert Quispe \\
\bottomrule
\end{tabularx}
}

\section{Control de la Calidad}
El control revisa el producto. Las pruebas automáticas usan \textbf{Jest} sobre el motor de puntaje; las manuales se ejecutan en \textbf{dos celulares Android reales}. Los casos se ejecutan en las semanas 7 (motor) y 15 (aceptación) y cada vez que cambia el motor.

\vspace{0.1cm}
{\footnotesize
\setlength{\tabcolsep}{3.5pt}
\rowcolors{2}{white}{GrisFila}
\noindent\begin{tabularx}{\textwidth}{@{}C{1.1cm} C{1.3cm} C{1.6cm} L@{}}
\toprule
\rowcolor{AzulClaro}
\textbf{Caso} & \textbf{Tipo} & \textbf{Requisito} & \textbf{Entrada y resultado esperado} \\
\midrule
CP-01 & \etiqueta{Rama2}{Jest} & RF-05 & Precio: S/ 140 da 100; S/ 850 da 0; S/ 495 da 50. \\
CP-02 & \etiqueta{Rama2}{Jest} & RF-05 & Cercanía: 50 m y 2 min dan 100; 3,000 m y 40 min dan 0. \\
CP-03 & \etiqueta{Rama2}{Jest} & RF-05 & Seguridad: nivel 1 da 20 y nivel 5 da 100. \\
CP-04 & \etiqueta{Rama2}{Jest} & RF-05 & Servicios: agua y luz sin internet dan 70; baño privado da 100 y compartido 0. \\
CP-05 & \etiqueta{Rama2}{Jest} & RF-06 & Bonificaciones y penalizaciones: ventana con luz suma 8; llegada hasta las 20:00 resta 20; condiciones no declaradas restan 25; el puntaje queda entre 0 y 100. \\
CP-06 & \etiqueta{Rama2}{Jest} & RF-04 & Pesos: si un criterio tiene peso 0, no influye en el puntaje. \\
CP-07 & \etiqueta{Rama2}{Jest} & RF-06 & Ranking: 5 cuartos de prueba quedan en el orden esperado. \\
CP-08 & \etiqueta{Rama4}{Manual} & RF-01 & Un arrendatario y un propietario se registran e inician sesión; cada uno ve solo las opciones de su rol. \\
CP-09 & \etiqueta{Rama4}{Manual} & RF-02 & Sin una foto o dato obligatorio el cuarto no se publica; con todos los datos se publica en 10 minutos o menos. \\
CP-10 & \etiqueta{Rama4}{Manual} & RF-03 & La distancia calculada al punto de referencia difiere en menos de 50 m de la medida en el mapa. \\
CP-11 & \etiqueta{Rama4}{Manual} & RF-07 & El detalle muestra el puntaje de los 5 criterios y cada bonificación y penalización aplicada. \\
CP-12 & \etiqueta{Rama4}{Manual} & RNF-01 & Con 500 cuartos de prueba, el ranking aparece en 250 ms o menos. \\
CP-13 & \etiqueta{Rama4}{Manual} & RNF-03 & Un propietario no puede editar ni ocultar cuartos de otro propietario. \\
CP-14 & \etiqueta{Rama4}{Manual} & RNF-02 & Todas las pantallas funcionan en un celular con Android 10 y en otro con una versión reciente. \\
\bottomrule
\end{tabularx}
}

\vspace{0.2cm}
\Needspace{9\baselineskip}
\noindent\textbf{Clasificación de defectos}

\vspace{0.1cm}
{\small
\noindent\begin{tabularx}{\textwidth}{@{}P{2.2cm} L P{3.3cm}@{}}
\toprule
\rowcolor{AzulClaro}
\textbf{Severidad} & \textbf{Descripción} & \textbf{Plazo de corrección} \\
\midrule
\etiqueta{SevAlta}{CRÍTICO} & Bloquea una función obligatoria, pierde datos o calcula mal el puntaje. & Antes del siguiente sprint \\
\etiqueta{SevMedia}{MAYOR} & Una función falla, pero existe una forma de continuar. & Dentro de 2 sprints \\
\etiqueta{SevBaja}{MENOR} & Defecto visual o de texto. & Antes de la entrega final \\
\bottomrule
\end{tabularx}
}

\section{Criterios de Aceptación}
\begin{itemize}[leftmargin=*]
    \item \textbf{MVP de la Unidad I (semana 8):} CP-01 a CP-07 aprobados y ranking en 250 ms o menos.
    \item \textbf{MVP final (semana 16):} los 14 casos aprobados, 0 defectos críticos y 3 o menos defectos mayores abiertos.
    \item \textbf{Piloto (semana 13):} 80\,\% o más de los 10 arrendatarios y de los 5 propietarios califica con 4 o 5 sobre 5.
    \item \textbf{Documentos:} revisados por otro integrante, sin errores de cifras y firmados por el equipo.
\end{itemize}

\Needspace{10\baselineskip}
\section{Definición de Terminado}
Un Issue se cierra solo cuando:
\begin{enumerate}[leftmargin=*]
    \item El código compila sin errores de TypeScript y sus pruebas de Jest pasan.
    \item Otro integrante aprobó el Pull Request.
    \item La función se probó en un celular Android real.
    \item El documento o la pantalla relacionada quedó actualizado en el repositorio.
\end{enumerate}

\Needspace{12\baselineskip}
\section{Costos de la Calidad}
Las horas forman parte del presupuesto aprobado (tarifa promedio de S/ 12.0625 por hora); no son un costo adicional. La meta es invertir más en prevención y evaluación para reducir las fallas.

\vspace{0.1cm}
{\small
\rowcolors{2}{white}{GrisFila}
\noindent\begin{tabularx}{\textwidth}{@{}P{2.6cm} L Q{1.1cm} Q{1.8cm} P{3.2cm}@{}}
\toprule
\rowcolor{AzulClaro}
\textbf{Categoría} & \textbf{Actividades} & \textbf{Horas} & \textbf{Costo S/} & \textbf{Proporción} \\
\midrule
Prevención & Elaboración de estándares y de este plan; práctica con Jest y Supabase & 16 & 193.00 & \tikz{\fill[Rama4] (0,0) rectangle (0.73,0.25);} 18\,\% \\
Evaluación & Pruebas del motor, pruebas de aceptación y revisión de código en cada Pull Request & 48 & 579.00 & \tikz{\fill[Rama2] (0,0) rectangle (2.20,0.25);} 55\,\% \\
Fallas internas & Corrección de errores encontrados antes del piloto & 16 & 193.00 & \tikz{\fill[SevMedia] (0,0) rectangle (0.73,0.25);} 18\,\% \\
Fallas externas & Correcciones por observaciones del piloto & 8 & 96.50 & \tikz{\fill[SevAlta] (0,0) rectangle (0.37,0.25);} 9\,\% \\
\midrule
\rowcolor{white}\textbf{Total} & & \textbf{88} & \textbf{1,061.50} & \\
\bottomrule
\end{tabularx}
}

\Needspace{12\baselineskip}
\section{Métricas de Seguimiento}
{\small
\rowcolors{2}{white}{GrisFila}
\noindent\begin{tabularx}{\textwidth}{@{}L C{2.4cm} P{2.8cm}@{}}
\toprule
\rowcolor{AzulClaro}
\textbf{Métrica} & \textbf{Meta} & \textbf{Frecuencia} \\
\midrule
Casos de prueba aprobados & 100\,\% en la entrega & Semanas 7, 15 y cada cambio \\
Defectos críticos abiertos & 0 & Martes \\
Cobertura de Jest en el motor de puntaje & 80\,\% o más & Cada Pull Request \\
Pull Requests revisados por otro integrante & 100\,\% & Martes \\
Satisfacción de los clientes en el piloto & 80\,\% o más & Semana 13 \\
\bottomrule
\end{tabularx}
}

\Needspace{24\baselineskip}
\section{Firmas}
\noindent En señal de conformidad con el presente documento, los integrantes del equipo firman a continuación:

\vspace{1.5cm}
\begin{center}
\begin{tabular}{c c c}
    \parbox{6.5cm}{\centering \rule{6.0cm}{0.6pt}\\[5pt] \textbf{Clyde Neil Paricahua Pari}\\[2pt] {\footnotesize Project Manager}}
    & \hspace{0.8cm} &
    \parbox{6.5cm}{\centering \rule{6.0cm}{0.6pt}\\[5pt] \textbf{Ruth Karina Apaza Solis}\\[2pt] {\footnotesize Analista de Planificación}}
    \\[2.0cm]
    \parbox{6.5cm}{\centering \rule{6.0cm}{0.6pt}\\[5pt] \textbf{Eliseo Tarqui Ajahuana}\\[2pt] {\footnotesize Especialista Técnico / Arquitectura}}
    & \hspace{0.8cm} &
    \parbox{6.5cm}{\centering \rule{6.0cm}{0.6pt}\\[5pt] \textbf{Viviana Mary Quisbert Quispe}\\[2pt] {\footnotesize Gestora de Comunicación, Riesgos y Documentación}}
\end{tabular}
\end{center}
\end{document}
